\datos{color}{}{}
\section{\textit{iRobot 510 PackBot}}
\textit{iRobot 510 PackBot}, desarrollado por \textit{iRobot Corporation} y presentado en 2007, Figura~\ref{packbot}, es un robot terrestre con locomoción por cadenas diseñado para operaciones de desactivación de explosivos, inspección, reconocimiento y tareas militares. Su robustez, modularidad y capacidad para operar en entornos hostiles lo han convertido en uno de los robots de intervención más utilizados en el ámbito militar y de protección civil. El sistema integra cámaras, sensores inerciales, sensores de proximidad y un ordenador central encargado de la navegación y el control. \\

Sus usos principales incluyen la inspección de artefactos explosivos, la exploración de edificios, la intervención en escenarios peligrosos y la asistencia en operaciones militares. PackBot ha sido desplegado en miles de misiones reales y es una de las plataformas más reconocidas en robótica de intervención. \\

A diferencia de robots distribuidos como \textit{RHex}, \textit{Crazyflie} o \textit{Aqua}, \textit{PackBot} implementa una arquitectura centralizada. El ordenador principal gestiona la percepción, la navegación y el control, siguiendo la descomposición vertical clásica. Los módulos de bajo nivel actúan como controladores esclavos que ejecutan órdenes del sistema central, sin autonomía local significativa. Esta arquitectura prioriza la fiabilidad, la seguridad y la supervisión humana, características esenciales en robots de intervención. \\

\begin{wrapfigure}{r}{0.34\textwidth}
\centering
\vspace{-10pt}
\includegraphics[width=0.34\textwidth]{Images/packbot.jpg}
\caption{\label{packbot}\textit{iRobot 510 PackBot}.}
\vspace{-10pt}
\end{wrapfigure}

\textit{PackBot} utiliza un sistema de locomoción basado en cadenas y ruedas auxiliares que le permiten superar obstáculos, subir escaleras y desplazarse por terrenos irregulares. Su diseño robusto y su arquitectura centralizada lo convierten en una herramienta fiable para operaciones críticas. \\

\subsection*{Referencias}
\cite{PackBot2007} \fullcite{PackBot2007}.\\
\cite{PackBot2012} \fullcite{PackBot2012}.